\begin{figure}[htbp]
\centering
\begin{tikzpicture}[
  stage/.style={
    draw,
    rounded corners,
    align=left,
    text width=14cm,
    inner sep=7pt,
    font=\small
  },
  >=Stealth
]

\node[stage] (national) {
  \textbf{1. National analysis}\\[3pt]
  \textbf{H1}: contemporaneous associations
  (Table~\ref{tab:annual}).\\
  592 country-years in 60 countries.\\[3pt]
  \textbf{H2}: longitudinal connectedness pathway
  (Table~\ref{tab:country-sem}).\\
  460 country-years in 60 countries.
};

\node[stage, below=6mm of national] (regional) {
  \textbf{2. Regional extension}\\[3pt]
  \textbf{H1-Reg}: contemporaneous associations
  (Table~\ref{tab:annual}).\\
  7,840 region-years in 953 regions and 59 countries.\\
  Sensitivity: replace national trends and common year effects with
  country-by-year effects, using the same sample
  (Table~\ref{tab:country-year-robustness}).\\[3pt]
  \textbf{H2-Reg}: longitudinal connectedness pathway
  (Table~\ref{tab:region-sem}).\\
  4,714 region-years in 796 regions and 59 countries.
};

\node[stage, below=6mm of regional] (averages) {
  \textbf{3. Comparison of regional period averages}\\[3pt]
  \textbf{H3}: associations between regional averages within countries
  (Table~\ref{tab:static}).\\
  682 regions in 42 countries, averaging observations over 2013--2022.\\
  Complementary analysis: static connectedness decomposition
  (Table~\ref{tab:static-sem}).
};

\node[stage, below=6mm of averages] (contrasts) {
  \textbf{4. Direct comparisons across well-being outcomes}\\[3pt]
  \textbf{H4}: compare the religious-expression coefficient for
  \emph{purpose} with those for \emph{happiness} and
  \emph{life satisfaction}
  (Table~\ref{tab:contrasts}).\\
  Contrasts are estimated separately for the national annual,
  regional annual, and regional-average association models.
  Regional contrasts are also repeated with country-by-year effects.
};

\draw[->] (national) -- (regional);
\draw[->] (regional) -- (averages);
\draw[->] (averages) -- (contrasts);

\end{tikzpicture}
\caption{Overview of the estimation strategy.
Arrows indicate the order of presentation.
All analyses use observations from 2013--2022.
The models are estimated separately, with sample sizes determined
by their respective inclusion criteria.}
\label{fig:estimation-strategy}
\end{figure}
